%%%%%%%%%%%%%%%%%%%%%%%%%%%%%%%%%%%%%%%%%%%%%%%%%%%
%%  复旦大学（硕士）学位论文模版 (UTF-8/XeLaTex)
%%  2013-03 richard@bbs.fudan.edu.cn
%%%%%%%%%%%%%%%%%%%%%%%%%%%%%%%%%%%%%%%%%%%%%%%%%%%

\documentclass[12pt,a4paper,oneside,openany]{book}

\usepackage{fontspec,xunicode,xltxtra}
\usepackage[slantfont,boldfont,CJKnumber]{xeCJK}
\usepackage[bookmarksnumbered]{hyperref}
\usepackage{graphicx}
\usepackage{titlesec} %% 用来产生中文章节样式
\usepackage{titletoc,tocbibind} %% MARK 增加两个package来处理目录
%\usepackage[titles]{tocloft} %% MARK 用来产生目录中的中文章节样式
\usepackage[numbers,sort&compress]{natbib} %% 根据使用gb7714.bst的要求添加
\usepackage{pdfpages} %% 用来嵌入封面和封底的pdf

\defaultfontfeatures{Mapping=tex-text} %% 是tex命令忽略双引号、连字符等
\setromanfont{Times New Roman}
\setCJKmainfont{宋体}
\settowidth{\parindent}{空格} %% 设置段首空格两个中文字符
\setCJKfamilyfont{youyuan}{黑体} %% 用来产生封面论文标题

%% 以下将水平方向左右边距全部调整为3cm
\setlength{\oddsidemargin}{0pt}
\setlength{\evensidemargin}{0pt}
\setlength{\marginparsep}{0pt}
\setlength{\hoffset}{0.46cm}
\setlength{\textwidth}{15cm}

%% 以下将垂直方向上下边距全部调整为3cm
\setlength{\voffset}{0.46cm}
\setlength{\topmargin}{0pt}
\setlength{\headheight}{0pt}
\setlength{\headsep}{0pt}
\setlength{\textheight}{23.7cm}

\linespread{1.3} %% 将行距调整至大约20磅的效果
\pagestyle{plain} %% 去掉页眉

%% 调整中文标题样式
\renewcommand{\contentsname}{目\hspace{1.5em}录}
\renewcommand{\chaptername}{第\CJKnumber{\thechapter}章}
\titleformat{\chapter}[block]{\bfseries\huge\filcenter}{\chaptertitlename}{1em}{}
\renewcommand{\bibname}{参考文献}
%% MARK 下面是利用tocloft的，删除掉
%\renewcommand{\cftchappresnum}{第\CJKnumber}
%\renewcommand{\cftchapaftersnum}{章}
%\renewcommand{\cftchapnumwidth}{4em}
\renewcommand{\figurename}{图}
\renewcommand{\tablename}{表}
\renewcommand{\appendixname}{附录\thechapter} %% MARK 附录格式

%% MARK 重新使用titletoc调整目录样式
\titlecontents{chapter}[0pt]{}{\contentspush{%
\ifnum0<0\thecontentslabel\relax
第\CJKnumber{\thecontentslabel}章\ %
\else
附录\thecontentslabel\ 
\fi
}}{}{\titlerule*[8pt]{.}\contentspage}


%% 用来产生封面的命令
%% 需要fudan.png图片文件
\newcommand{\makefrontcover}[8]{\thispagestyle{empty} %% 去掉页码
\begin{flushright} 
\fontsize{9pt}{11pt}\selectfont
\begin{tabular}{rl}
学校代码：& #1 \\
学\hspace{2em}号：& #2 \\
\end{tabular}
\end{flushright}
\vspace{1.2cm}
\begin{center}
\includegraphics[width=7.68cm,height=2.86cm]{fudan.png}
\end{center}
\vspace{0.3cm}
\begin{center}
\fontsize{22pt}{22pt}\selectfont
硕\hspace{.5em}士\hspace{.5em}学\hspace{.5em}位\hspace{.5em}论\hspace{.5em}文
\end{center}
\vspace{1.5cm}
\begin{center}
\CJKfamily{youyuan}
\fontsize{18pt}{18pt}\selectfont
\textbf{#3}
\end{center}
\vspace{4.1cm}
\begin{flushright}
\fontsize{14pt}{24pt}\selectfont
\begin{tabular}{p{90pt}p{200pt}}
院\hspace{3.5em}系：& #4 \\
专\hspace{3.5em}业：& #5 \\
姓\hspace{3.5em}名：& #6 \\
指\hspace{.5em}导\hspace{.5em}教\hspace{.5em}师：& #7 \\
完\hspace{.5em}成\hspace{.5em}日\hspace{.5em}期：& #8 \\
\end{tabular}
\end{flushright}\newpage}

%% 引用gb7714标准的参考文献格式
%% bst文件下载自：http://bbs.ctex.org/forum.php?mod=viewthread&tid=33571
%% 进行了部分修改以减少编译错误
\bibliographystyle{gb7714}

\title{论文标题}
\author{
指导小组成员名单\vspace{1em}\\
某\hspace{1em}人 \hspace{1em} 教\hspace{1em}授\\
某某人 \hspace{1em} 副教授\\
}
\date{}

\begin{document}

%% 嵌入封面frontcover.pdf
%% \includepdf[noautoscale]{frontcover}
%% 也可以选择使用下面的命令产生封面，需要fudan.png文件
\makefrontcover{10246}{1020124XXXX}%
{论文标题}%
{计算机科学技术学院}%
{计算机软件与理论}%
{某学生}%
{某某人\hspace{1.5em}教授}%
{2013年4月22日}

%% 扉页
\maketitle

%% 目录
\tableofcontents

%% 去掉摘要章节的章节编号
\frontmatter

%% 中文摘要
\chapter{摘\hspace{1.5em}要}

\textbf{【关键词】}

\textbf{【中图分类号】} TP183

%% 英文摘要
\chapter{Abstract}

\textbf{Keywords:}

\textbf{CLC Number:} TP183

%% 正文从这里开始
\mainmatter

\chapter{章的标题}

\section{节的标题}

\subsection{小节的标题}

内容。

%% 一个表格
\begin{table}[!h]
\begin{center}
\begin{tabular}{r|c}
\hline\hline
方法 & 结果 \\
\hline
他们的 & 还可以 \\
你们的 & 还可以 \\
我们的 & 非常非常的好啊 \\
\hline\hline
\end{tabular}
\end{center}
\caption{实验结果}
\end{table}

%% 一张图片
\begin{figure}[!h]
\begin{center}
\fbox{\rule{0pt}{2in} \rule{0.9\linewidth}{0pt}}
\end{center}
\caption{一个图示}
\end{figure}

\nocite{*} %% 列出全部（包括未引用的）参考文献

\begin{appendix}

\chapter{附录}

\end{appendix}

%% 去掉附录部分的章节编号
\backmatter

%% 将参考文献列入目录 (MARK 注释掉了，不在需要)
%% MARK \addcontentsline{toc}{chapter}{参考文献}
\bibliography{ref}

\chapter{致\hspace{1.5em}谢}

%% 嵌入封底backcover.pdf
\includepdf[noautoscale]{backcover}

\end{document}
